
\documentclass{article}
\usepackage{colortbl}
\usepackage{makecell}
\usepackage{multirow}
\usepackage{supertabular}

\begin{document}

\newcounter{utterance}

\centering \large Interaction Transcript for game `cladder', experiment `full\_v1.5\_default', episode 7275 with Qwen3.5{-}9B.
\vspace{24pt}

{ \footnotesize  \setcounter{utterance}{1}
\setlength{\tabcolsep}{0pt}
\begin{supertabular}{c@{$\;$}|p{.15\linewidth}@{}p{.15\linewidth}p{.15\linewidth}p{.15\linewidth}p{.15\linewidth}p{.15\linewidth}}
    \# & \multicolumn{2}{c}{Player} && \multicolumn{2}{c}{Game Master} \\
    \hline

    \theutterance \stepcounter{utterance}  
    & & & \multicolumn{4}{p{0.6\linewidth}}{
        \cellcolor[rgb]{0.9,0.9,0.9}{
            \makecell[{{p{\linewidth}}}]{
                \texttt{\tiny{[P1$\langle$GM]}}
                \texttt{You are an expert in causal inference. The following question is not a typical commonsense query, but rather a meticulously designed question created by a professor specializing in causal inference, intended to assess the students' mastery of the course content.} \\
\texttt{Imagine a self{-}contained, hypothetical world with only the following conditions, and without any unmentioned factors or causal relationships: Yupt has a direct effect on zupj and muvq. Zupj has a direct effect on muvq. The overall probability of yupt is 1\%. For those who are not yupt, the probability of muvq is 33\%. For those who are yupt, the probability of muvq is 59\%. Is muvq more likely than not muvq overall?} \\
\texttt{Start your answer with 'Yes' or 'No', followed by additional reasoning or evidence to support your explanation.} \\
            }
        }
    }
     \\ \\

    \theutterance \stepcounter{utterance}  
    & \multicolumn{4}{p{0.6\linewidth}}{
        \cellcolor[rgb]{0.561,0.737,0.561}{
            \makecell[{{p{\linewidth}}}]{
                \texttt{\tiny{[P1$\rangle$GM]}}
                \texttt{No} \\
\\ 
\texttt{**Reasoning:**} \\
\\ 
\texttt{To determine if \$muvq\$ is more likely than not \$muvq\$ overall, we must calculate the total probability \$P(muvq)\$ using the Law of Total Probability. This law states that the total probability of an event is the sum of the probabilities of that event occurring in each mutually exclusive partition of the sample space, weighted by the probability of each partition.} \\
\\ 
\texttt{The problem provides two mutually exclusive and exhaustive conditions for the variable \$yupt\$:} \\
\texttt{1.  **Condition A:** The individual is \$yupt\$.} \\
\texttt{    *   \$P(yupt) = 1\textbackslash{}\% = 0.01\$} \\
\texttt{    *   \$P(muvq | yupt) = 59\textbackslash{}\% = 0.59\$} \\
\texttt{2.  **Condition B:** The individual is not \$yupt\$ (denoted as \$\textbackslash{}neg yupt\$).} \\
\texttt{    *   \$P(\textbackslash{}neg yupt) = 1 {-} P(yupt) = 1 {-} 0.01 = 0.99\$} \\
\texttt{    *   \$P(muvq | \textbackslash{}neg yupt) = 33\textbackslash{}\% = 0.33\$} \\
\\ 
\texttt{Now, we calculate the overall probability \$P(muvq)\$:} \\
\texttt{\$\$P(muvq) = P(muvq | yupt) \textbackslash{}times P(yupt) + P(muvq | \textbackslash{}neg yupt)} \\
            }
        }
    }
    & & \\ \\

    \theutterance \stepcounter{utterance}  
    & & & \multicolumn{2}{p{0.3\linewidth}}{
        \cellcolor[rgb]{0.9,0.9,0.9}{
            \makecell[{{p{\linewidth}}}]{
                \texttt{\tiny{[GM$|$GM]}}
                \texttt{No} \\
            }
        }
    }
    & & \\ \\

    \theutterance \stepcounter{utterance}  
    & & & \multicolumn{2}{p{0.3\linewidth}}{
        \cellcolor[rgb]{0.9,0.9,0.9}{
            \makecell[{{p{\linewidth}}}]{
                \texttt{\tiny{[GM$|$GM]}}
                \texttt{no} \\
            }
        }
    }
    & & \\ \\

    \theutterance \stepcounter{utterance}  
    & & & \multicolumn{2}{p{0.3\linewidth}}{
        \cellcolor[rgb]{0.9,0.9,0.9}{
            \makecell[{{p{\linewidth}}}]{
                \texttt{\tiny{[GM$|$GM]}}
                \texttt{game\_result = WIN} \\
            }
        }
    }
    & & \\ \\

\end{supertabular}
}

\end{document}
